\chapter{Estado del arte}
\label{cap:estado_arte}

Este capítulo sitúa el trabajo en dos tradiciones que dialogan a lo largo de la memoria: la del
\emph{factor investing}, que aporta el listón de comparación y una explicación económica de qué
mueve los retornos, y la del \emph{aprendizaje automático aplicado a los mercados}, que aporta la
herramienta y, con ella, un catálogo de patologías metodológicas que este proyecto se propone
evitar. El capítulo termina justificando por qué se adopta el \rankic{} fuera de muestra, y no la
rentabilidad, como criterio de evidencia.

\section{Primas de factor como listón de comparación}

La idea de que el retorno esperado de una acción se explica por su exposición a un conjunto
reducido de \emph{factores} sistemáticos es uno de los resultados mejor establecidos de las
finanzas empíricas. El modelo de tres factores de \textcite{fama1993common} amplió el CAPM
clásico añadiendo al factor de mercado un factor de \emph{tamaño} ---las empresas pequeñas
tienden a rendir más que las grandes, un patrón ya documentado por \textcite{banz1981relationship}---
y un factor de \emph{valor}, capturado por el cociente entre valor contable y valor de mercado
\parencite{fama1992cross}. A esa base se sumó el factor de \emph{momentum}: las acciones que han
subido en los últimos meses tienden a seguir subiendo en el corto plazo
\parencite{jegadeesh1993returns}, un efecto que \textcite{carhart1997persistence} incorporó como
cuarto factor. Más recientemente, la literatura ha consolidado un factor de \emph{calidad}
---empresas rentables, estables y poco apalancadas--- como fuente de prima diferenciada
\parencite{asness2013quality}.

Estos factores son relevantes para el presente trabajo por dos motivos. Primero, definen el
\textbf{listón}: cualquier sistema de selección de acciones debe compararse con lo que se obtendría
explotando de forma mecánica estas primas conocidas, no solo con el índice. Segundo, agrupan de
forma natural los estilos de inversión ---calidad, valor, momentum y sus combinaciones, como la
estrategia \emph{growth at a reasonable price} o GARP--- que este proyecto encapsula en agentes
especializados (capítulo~\ref{cap:agentes}) y en perfiles de inversor
(capítulo~\ref{cap:diseno_experimental}).

\section{Aprendizaje automático aplicado a la selección de activos}

La disponibilidad de datos y de métodos no lineales ha impulsado una línea de investigación que
sustituye el modelo lineal de factores por algoritmos de aprendizaje automático capaces de captar
interacciones entre variables. El trabajo de referencia de \textcite{gu2020empirical} compara de
forma sistemática un amplio abanico de métodos ---desde regresiones penalizadas hasta árboles con
\emph{gradient boosting} y redes neuronales--- para predecir retornos en el corte transversal, y
documenta que los métodos no lineales, en particular los basados en árboles, tienden a superar a
los lineales cuando la relación entre factores y retorno no es aditiva. Entre esos métodos, el
\emph{gradient boosting} sobre árboles de decisión ---y su implementación eficiente LightGBM
\parencite{ke2017lightgbm}--- se ha popularizado por su capacidad de modelar interacciones con un
coste computacional moderado, razón por la que este proyecto lo adopta como modelo de sus agentes
(capítulo~\ref{cap:agentes}).

Ahora bien, la potencia de estos métodos convive con una fragilidad conocida: cuanto más flexible
es el modelo, más fácil es que ajuste ruido en lugar de señal, sobre todo cuando el número de
periodos verdaderamente independientes es pequeño. Esa tensión motiva la atención metodológica que
vertebra el resto del capítulo.

\section{Patologías metodológicas recurrentes}

La literatura crítica ha identificado un conjunto de errores que inflan sistemáticamente el
rendimiento aparente de una estrategia y que, si no se controlan, hacen irreproducibles los
resultados. Tres son especialmente relevantes para este trabajo:

\begin{itemize}
  \item \textbf{Sesgo de anticipación} (\emph{lookahead bias}). Usar en una fecha información que
  en realidad no era observable ese día ---por ejemplo, un dato fundamental antes de que la empresa
  lo publicara, o la composición actual de un índice para decidir en el pasado--- introduce una
  ventaja irreal. Es el sesgo que el diseño \emph{point-in-time} del capítulo~\ref{cap:metodologia}
  ataca de raíz.

  \item \textbf{Sesgo de supervivencia} (\emph{survivorship bias}). Estudiar solo las empresas que
  han sobrevivido hasta hoy sobreestima el rendimiento, porque excluye a las que quebraron o fueron
  excluidas del índice \parencite{brown1992survivorship}. Este proyecto lo mide por año en lugar de
  limitarse a declararlo (capítulo~\ref{cap:datos}).

  \item \textbf{Sobreajuste por selección} (\emph{backtest overfitting}). Probar muchas
  configuraciones y quedarse con la mejor produce, casi con seguridad, un resultado atractivo aunque
  no exista señal alguna: el máximo de muchos intentos aleatorios es alto por construcción.
  \textcite{lopezdeprado2014pseudo} formalizan cómo un número suficiente de pruebas garantiza un
  \emph{backtest} espectacular sin capacidad predictiva real, y \textcite{harvey2016cross}
  documentan cómo la multiplicidad de contrastes ha poblado la literatura de factores que no
  replican. La respuesta metodológica ---validación cruzada respetuosa con el tiempo, reserva de un
  periodo que no interviene en ninguna decisión y contrastes de significancia--- se recoge en
  \textcite{lopezdeprado2018advances} y se aplica en el capítulo~\ref{cap:diseno_experimental}.
\end{itemize}

\section{Por qué el \texorpdfstring{\rankic}{rank-IC} fuera de muestra como criterio de evidencia}

De las patologías anteriores se desprende la decisión de evidencia de este TFM. Reportar
únicamente la rentabilidad de un \emph{backtest} es frágil frente a las tres: un artefacto de datos,
una empresa desaparecida o una configuración afortunada pueden producir una curva de rentabilidad
convincente sin que el modelo haya aprendido nada. Por eso se adopta como métrica principal el
\rankic{} \emph{fuera de muestra}, que pregunta directamente si el modelo ordena bien los activos,
y se le exige superar dos filtros: significancia estadística mediante remuestreo por bloques ---que
respeta el solapamiento temporal de las cohortes--- y estabilidad a lo largo de distintos periodos.
La rentabilidad se conserva, pero como \emph{consecuencia} a reportar tras depurar artefactos, nunca
como criterio de selección. Esta postura ---separar lo que el sistema \emph{aprende} de lo que
\emph{rinde}--- es la que articula los capítulos de método y resultados que siguen.
